% --- Archon physics formalization source begin ---
% archon:physics
% archon:covers PhyXMiniProblems/problem_phyx_mini_0579.lean
% archon:source-report reports/phyx_mini/problem_phyx_mini_0579.source.json

\chapter{Physics problem phyx\_mini\_0579}
\label{ch:PhyXMiniProblems_problem_phyx_mini_0579}

\paragraph{Problem source.}
\#\# Physical scenario

Where sunlight shines on the atmosphere of Mars, carbon dioxide molecules at an altitude of about \$75 \textbackslash{}, \textbackslash{}text\{km\}\$ undergo natural laser action. The energy levels involved in the action are shown figure; population inversion occurs between energy levels \$E\_2\$ and \$E\_1\$.

\#\# Question

At what wave-length does lasing occur?

\#\# Answer choices

- A: A: 4.29μm

- B: B: 5.91μm

- C: C: 10.0μm

- D: D: 8.60μm

\#\# Figure caption (auxiliary; use the image as primary evidence)

The image is an energy level diagram with three horizontal lines representing energy levels. 

1. **Energy Levels**:
   - The bottom line is labeled \textbackslash{}(E\_0 = 0\textbackslash{}).
   - The middle line is labeled \textbackslash{}(E\_1 = 0.165 \textbackslash{}, \textbackslash{}text\{eV\}\textbackslash{}).
   - The top line is labeled \textbackslash{}(E\_2 = 0.289 \textbackslash{}, \textbackslash{}text\{eV\}\textbackslash{}).

2. **Arrows**:
   - An upward arrow extends from \textbackslash{}(E\_0\textbackslash{}) to \textbackslash{}(E\_2\textbackslash{}), indicating a transition from the lowest energy state to the highest.
   - A downward arrow extends from \textbackslash{}(E\_2\textbackslash{}) to \textbackslash{}(E\_1\textbackslash{}), indicating a transition from the highest energy level to the middle one.

These arrows suggest absorption and emission processes typically seen in quantum mechanical systems or electron transitions in atoms/molecules.

\#\# Dataset metadata

Quantum Phenomena; Physical Model Grounding Reasoning, Multi-Formula Reasoning

\paragraph{Recorded answer/context.}
C: C: 10.0μm

\paragraph{Figure/image path.}
/root/proposal\_for\_physic/science-mango/phyx\_mini\_run/phyx\_data/test\_image/579.png

\paragraph{Formalization target.}
create a compiling Lean file with sorry bodies at `PhyXMiniProblems/problem\_phyx\_mini\_0579.lean`. The Lean declarations must preserve the physical quantities, dimensions or dimensional roles, figure labels, governing-law hypotheses, and final relation expressed by this problem.
Use Mathlib/Physlib names found through LeanExplore where available. If a physics API is missing, introduce faithful local abstractions rather than scalar placeholder aliases.

\begin{theorem}[Physics formalization target]
\label{thm:physics:phyx_mini_0579:target}
\lean{PhyXMiniProblems.ProblemPhyXMini0579.problem_phyx_mini_0579}
\uses{def:physics:phyx-mini-0579:phyxminiproblems-problemphyxmini0579-lengthinmeters, def:physics:phyx-mini-0579:phyxminiproblems-problemphyxmini0579-electronvoltinjoules, def:physics:phyx-mini-0579:phyxminiproblems-problemphyxmini0579-energyinelectronvolts, def:physics:phyx-mini-0579:phyxminiproblems-problemphyxmini0579-planckconstantsivalue, def:physics:phyx-mini-0579:phyxminiproblems-problemphyxmini0579-speedoflightinmeterspersecond, def:physics:phyx-mini-0579:phyxminiproblems-problemphyxmini0579-marsco2lasersetup, def:physics:phyx-mini-0579:phyxminiproblems-problemphyxmini0579-matchesmarsco2laserscenario, def:physics:phyx-mini-0579:phyxminiproblems-problemphyxmini0579-matchessuppliedco2energylevelfigure, def:physics:phyx-mini-0579:phyxminiproblems-problemphyxmini0579-hasphysicalco2laserquantities, def:physics:phyx-mini-0579:phyxminiproblems-problemphyxmini0579-usesordinaryplanckconstant, def:physics:phyx-mini-0579:phyxminiproblems-problemphyxmini0579-satisfiesco2laserlaws, def:physics:phyx-mini-0579:phyxminiproblems-problemphyxmini0579-isuniquematchingdisplayedwavelength}
The assigned autoformalize agent should translate this physics problem into Lean declarations in the covered file, with theorem and lemma proof bodies written as `by sorry`.
\end{theorem}
\begin{proof}
This is an autoformalization task, not a proof task. Produce statements that can later be proved by the physics prover without weakening the physical model.
\end{proof}

% --- Archon declaration topology begin ---
\section{Declaration topology}
\begin{definition}[Length Quantity]
  \label{def:physics:phyx-mini-0579:phyxminiproblems-problemphyxmini0579-lengthquantity}
  \lean{PhyXMiniProblems.ProblemPhyXMini0579.LengthQuantity}
  A nonnegative physical length, used for altitude and wavelength
\end{definition}
\begin{proof}
  This is the declared model definition of Length Quantity.
\end{proof}
\begin{definition}[Frequency Quantity]
  \label{def:physics:phyx-mini-0579:phyxminiproblems-problemphyxmini0579-frequencyquantity}
  \lean{PhyXMiniProblems.ProblemPhyXMini0579.FrequencyQuantity}
  A nonnegative physical frequency, carrying inverse-time dimension
\end{definition}
\begin{proof}
  This is the declared model definition of Frequency Quantity.
\end{proof}
\begin{definition}[action Dimension]
  \label{def:physics:phyx-mini-0579:phyxminiproblems-problemphyxmini0579-actiondimension}
  \lean{PhyXMiniProblems.ProblemPhyXMini0579.actionDimension}
  The physical dimension of action, energy * time
\end{definition}
\begin{proof}
  This is the declared model definition of action Dimension.
\end{proof}
\begin{definition}[Planck Constant Quantity]
  \label{def:physics:phyx-mini-0579:phyxminiproblems-problemphyxmini0579-planckconstantquantity}
  \lean{PhyXMiniProblems.ProblemPhyXMini0579.PlanckConstantQuantity}
  \uses{def:physics:phyx-mini-0579:phyxminiproblems-problemphyxmini0579-actiondimension}
  A nonnegative physical quantity with the dimension of Planck's constant
\end{definition}
\begin{proof}
  This is the declared model definition of Planck Constant Quantity.
\end{proof}
\begin{definition}[length Readout]
  \label{def:physics:phyx-mini-0579:phyxminiproblems-problemphyxmini0579-lengthreadout}
  \lean{PhyXMiniProblems.ProblemPhyXMini0579.lengthReadout}
  \uses{def:physics:phyx-mini-0579:phyxminiproblems-problemphyxmini0579-lengthquantity}
  Read a physical length in a selected length unit
\end{definition}
\begin{proof}
  This is the declared model definition of length Readout.
\end{proof}
\begin{definition}[length In Meters]
  \label{def:physics:phyx-mini-0579:phyxminiproblems-problemphyxmini0579-lengthinmeters}
  \lean{PhyXMiniProblems.ProblemPhyXMini0579.lengthInMeters}
  \uses{def:physics:phyx-mini-0579:phyxminiproblems-problemphyxmini0579-lengthquantity, def:physics:phyx-mini-0579:phyxminiproblems-problemphyxmini0579-lengthreadout}
  Read a physical length in metres
\end{definition}
\begin{proof}
  This is the declared model definition of length In Meters.
\end{proof}
\begin{definition}[length In Kilometers]
  \label{def:physics:phyx-mini-0579:phyxminiproblems-problemphyxmini0579-lengthinkilometers}
  \lean{PhyXMiniProblems.ProblemPhyXMini0579.lengthInKilometers}
  \uses{def:physics:phyx-mini-0579:phyxminiproblems-problemphyxmini0579-lengthquantity, def:physics:phyx-mini-0579:phyxminiproblems-problemphyxmini0579-lengthreadout}
  Read a physical length in kilometres
\end{definition}
\begin{proof}
  This is the declared model definition of length In Kilometers.
\end{proof}
\begin{definition}[length In Micrometers]
  \label{def:physics:phyx-mini-0579:phyxminiproblems-problemphyxmini0579-lengthinmicrometers}
  \lean{PhyXMiniProblems.ProblemPhyXMini0579.lengthInMicrometers}
  \uses{def:physics:phyx-mini-0579:phyxminiproblems-problemphyxmini0579-lengthquantity, def:physics:phyx-mini-0579:phyxminiproblems-problemphyxmini0579-lengthreadout}
  Read a physical wavelength in micrometres
\end{definition}
\begin{proof}
  This is the declared model definition of length In Micrometers.
\end{proof}
\begin{definition}[energy In Joules]
  \label{def:physics:phyx-mini-0579:phyxminiproblems-problemphyxmini0579-energyinjoules}
  \lean{PhyXMiniProblems.ProblemPhyXMini0579.energyInJoules}
  Read a physical energy in coherent-SI joules
\end{definition}
\begin{proof}
  This is the declared model definition of energy In Joules.
\end{proof}
\begin{definition}[electron Volt In Joules]
  \label{def:physics:phyx-mini-0579:phyxminiproblems-problemphyxmini0579-electronvoltinjoules}
  \lean{PhyXMiniProblems.ProblemPhyXMini0579.electronVoltInJoules}
  \uses{def:physics:phyx-mini-0579:phyxminiproblems-problemphyxmini0579-energyinjoules}
  The joule readout of Physlib's exact dimensionful electron volt
\end{definition}
\begin{proof}
  This is the declared model definition of electron Volt In Joules.
\end{proof}
\begin{definition}[energy In Electron Volts]
  \label{def:physics:phyx-mini-0579:phyxminiproblems-problemphyxmini0579-energyinelectronvolts}
  \lean{PhyXMiniProblems.ProblemPhyXMini0579.energyInElectronVolts}
  \uses{def:physics:phyx-mini-0579:phyxminiproblems-problemphyxmini0579-energyinjoules, def:physics:phyx-mini-0579:phyxminiproblems-problemphyxmini0579-electronvoltinjoules}
  Read a physical energy in electron volts
\end{definition}
\begin{proof}
  This is the declared model definition of energy In Electron Volts.
\end{proof}
\begin{definition}[frequency In Hertz]
  \label{def:physics:phyx-mini-0579:phyxminiproblems-problemphyxmini0579-frequencyinhertz}
  \lean{PhyXMiniProblems.ProblemPhyXMini0579.frequencyInHertz}
  \uses{def:physics:phyx-mini-0579:phyxminiproblems-problemphyxmini0579-frequencyquantity}
  Read a physical frequency in hertz
\end{definition}
\begin{proof}
  This is the declared model definition of frequency In Hertz.
\end{proof}
\begin{definition}[planck Constant In Joule Seconds]
  \label{def:physics:phyx-mini-0579:phyxminiproblems-problemphyxmini0579-planckconstantinjouleseconds}
  \lean{PhyXMiniProblems.ProblemPhyXMini0579.planckConstantInJouleSeconds}
  \uses{def:physics:phyx-mini-0579:phyxminiproblems-problemphyxmini0579-planckconstantquantity}
  Read a physical action in joule-seconds
\end{definition}
\begin{proof}
  This is the declared model definition of planck Constant In Joule Seconds.
\end{proof}
\begin{definition}[planck Constant SIValue]
  \label{def:physics:phyx-mini-0579:phyxminiproblems-problemphyxmini0579-planckconstantsivalue}
  \lean{PhyXMiniProblems.ProblemPhyXMini0579.planckConstantSIValue}
  The exact SI value assigned to the ordinary Planck constant
\end{definition}
\begin{proof}
  This is the declared model definition of planck Constant SIValue.
\end{proof}
\begin{definition}[speed Of Light In Meters Per Second]
  \label{def:physics:phyx-mini-0579:phyxminiproblems-problemphyxmini0579-speedoflightinmeterspersecond}
  \lean{PhyXMiniProblems.ProblemPhyXMini0579.speedOfLightInMetersPerSecond}
  Physlib's exact speed of light, read in metres per second
\end{definition}
\begin{proof}
  This is the declared model definition of speed Of Light In Meters Per Second.
\end{proof}
\begin{definition}[Planet]
  \label{def:physics:phyx-mini-0579:phyxminiproblems-problemphyxmini0579-planet}
  \lean{PhyXMiniProblems.ProblemPhyXMini0579.Planet}
  Planet on which the illuminated atmospheric molecules are located
\end{definition}
\begin{proof}
  This is the declared model definition of Planet.
\end{proof}
\begin{definition}[Molecular Environment]
  \label{def:physics:phyx-mini-0579:phyxminiproblems-problemphyxmini0579-molecularenvironment}
  \lean{PhyXMiniProblems.ProblemPhyXMini0579.MolecularEnvironment}
  Environment occupied by the illuminated molecules
\end{definition}
\begin{proof}
  This is the declared model definition of Molecular Environment.
\end{proof}
\begin{definition}[Molecule Species]
  \label{def:physics:phyx-mini-0579:phyxminiproblems-problemphyxmini0579-moleculespecies}
  \lean{PhyXMiniProblems.ProblemPhyXMini0579.MoleculeSpecies}
  Molecular species undergoing the natural laser action
\end{definition}
\begin{proof}
  This is the declared model definition of Molecule Species.
\end{proof}
\begin{definition}[Pump Source]
  \label{def:physics:phyx-mini-0579:phyxminiproblems-problemphyxmini0579-pumpsource}
  \lean{PhyXMiniProblems.ProblemPhyXMini0579.PumpSource}
  External mechanism that pumps the molecular energy levels
\end{definition}
\begin{proof}
  This is the declared model definition of Pump Source.
\end{proof}
\begin{definition}[Molecular Energy Level]
  \label{def:physics:phyx-mini-0579:phyxminiproblems-problemphyxmini0579-molecularenergylevel}
  \lean{PhyXMiniProblems.ProblemPhyXMini0579.MolecularEnergyLevel}
  The three horizontal molecular energy levels in the supplied diagram
\end{definition}
\begin{proof}
  This is the declared model definition of Molecular Energy Level.
\end{proof}
\begin{definition}[Figure Transition Arrow]
  \label{def:physics:phyx-mini-0579:phyxminiproblems-problemphyxmini0579-figuretransitionarrow}
  \lean{PhyXMiniProblems.ProblemPhyXMini0579.FigureTransitionArrow}
  The two arrows visible in the supplied energy-level diagram
\end{definition}
\begin{proof}
  This is the declared model definition of Figure Transition Arrow.
\end{proof}
\begin{definition}[Arrow Direction]
  \label{def:physics:phyx-mini-0579:phyxminiproblems-problemphyxmini0579-arrowdirection}
  \lean{PhyXMiniProblems.ProblemPhyXMini0579.ArrowDirection}
  Vertical direction in which an arrow is drawn
\end{definition}
\begin{proof}
  This is the declared model definition of Arrow Direction.
\end{proof}
\begin{definition}[expected Arrow Initial Level]
  \label{def:physics:phyx-mini-0579:phyxminiproblems-problemphyxmini0579-expectedarrowinitiallevel}
  \lean{PhyXMiniProblems.ProblemPhyXMini0579.expectedArrowInitialLevel}
  \uses{def:physics:phyx-mini-0579:phyxminiproblems-problemphyxmini0579-molecularenergylevel, def:physics:phyx-mini-0579:phyxminiproblems-problemphyxmini0579-figuretransitionarrow}
  The initial level of each transition represented in the figure
\end{definition}
\begin{proof}
  This is the declared model definition of expected Arrow Initial Level.
\end{proof}
\begin{definition}[expected Arrow Final Level]
  \label{def:physics:phyx-mini-0579:phyxminiproblems-problemphyxmini0579-expectedarrowfinallevel}
  \lean{PhyXMiniProblems.ProblemPhyXMini0579.expectedArrowFinalLevel}
  \uses{def:physics:phyx-mini-0579:phyxminiproblems-problemphyxmini0579-molecularenergylevel, def:physics:phyx-mini-0579:phyxminiproblems-problemphyxmini0579-figuretransitionarrow}
  The final level of each transition represented in the figure
\end{definition}
\begin{proof}
  This is the declared model definition of expected Arrow Final Level.
\end{proof}
\begin{definition}[expected Arrow Direction]
  \label{def:physics:phyx-mini-0579:phyxminiproblems-problemphyxmini0579-expectedarrowdirection}
  \lean{PhyXMiniProblems.ProblemPhyXMini0579.expectedArrowDirection}
  \uses{def:physics:phyx-mini-0579:phyxminiproblems-problemphyxmini0579-figuretransitionarrow, def:physics:phyx-mini-0579:phyxminiproblems-problemphyxmini0579-arrowdirection}
  The direction in which each transition arrow is drawn
\end{definition}
\begin{proof}
  This is the declared model definition of expected Arrow Direction.
\end{proof}
\begin{definition}[CO2 Energy Level Figure]
  \label{def:physics:phyx-mini-0579:phyxminiproblems-problemphyxmini0579-co2energylevelfigure}
  \lean{PhyXMiniProblems.ProblemPhyXMini0579.CO2EnergyLevelFigure}
  \uses{def:physics:phyx-mini-0579:phyxminiproblems-problemphyxmini0579-molecularenergylevel, def:physics:phyx-mini-0579:phyxminiproblems-problemphyxmini0579-figuretransitionarrow, def:physics:phyx-mini-0579:phyxminiproblems-problemphyxmini0579-arrowdirection}
  Literal data carried by image 579. The printed energies are scalar electron-volt readouts, separate from the physical energies in the setup
\end{definition}
\begin{proof}
  This is the declared model definition of CO2 Energy Level Figure.
\end{proof}
\begin{definition}[Mars CO2 Laser Setup]
  \label{def:physics:phyx-mini-0579:phyxminiproblems-problemphyxmini0579-marsco2lasersetup}
  \lean{PhyXMiniProblems.ProblemPhyXMini0579.MarsCO2LaserSetup}
  \uses{def:physics:phyx-mini-0579:phyxminiproblems-problemphyxmini0579-lengthquantity, def:physics:phyx-mini-0579:phyxminiproblems-problemphyxmini0579-frequencyquantity, def:physics:phyx-mini-0579:phyxminiproblems-problemphyxmini0579-planckconstantquantity, def:physics:phyx-mini-0579:phyxminiproblems-problemphyxmini0579-planet, def:physics:phyx-mini-0579:phyxminiproblems-problemphyxmini0579-molecularenvironment, def:physics:phyx-mini-0579:phyxminiproblems-problemphyxmini0579-moleculespecies, def:physics:phyx-mini-0579:phyxminiproblems-problemphyxmini0579-pumpsource, def:physics:phyx-mini-0579:phyxminiproblems-problemphyxmini0579-molecularenergylevel, def:physics:phyx-mini-0579:phyxminiproblems-problemphyxmini0579-co2energylevelfigure}
  The wavelength, frequency, and photon energy of the lasing radiation are independent observables. The hypotheses below relate them through physical laws; no answer-choice wavelength is stored in this structure
\end{definition}
\begin{proof}
  This is the declared model definition of Mars CO2 Laser Setup.
\end{proof}
\begin{definition}[Is About]
  \label{def:physics:phyx-mini-0579:phyxminiproblems-problemphyxmini0579-isabout}
  \lean{PhyXMiniProblems.ProblemPhyXMini0579.IsAbout}
  A scalar readout lies within a stated nonnegative tolerance of a value
\end{definition}
\begin{proof}
  This is the declared model definition of Is About.
\end{proof}
\begin{definition}[Matches Mars CO2 Laser Scenario]
  \label{def:physics:phyx-mini-0579:phyxminiproblems-problemphyxmini0579-matchesmarsco2laserscenario}
  \lean{PhyXMiniProblems.ProblemPhyXMini0579.MatchesMarsCO2LaserScenario}
  \uses{def:physics:phyx-mini-0579:phyxminiproblems-problemphyxmini0579-lengthinkilometers, def:physics:phyx-mini-0579:phyxminiproblems-problemphyxmini0579-marsco2lasersetup, def:physics:phyx-mini-0579:phyxminiproblems-problemphyxmini0579-isabout}
  The prose data: illuminated carbon dioxide in the Martian atmosphere near 75 km, with a population inversion from E₁ to E₂ and natural laser action present. These are scenario assumptions, not conclusions about the wavelength
\end{definition}
\begin{proof}
  This is the declared model definition of Matches Mars CO2 Laser Scenario.
\end{proof}
\begin{definition}[Matches Supplied CO2 Energy Level Figure]
  \label{def:physics:phyx-mini-0579:phyxminiproblems-problemphyxmini0579-matchessuppliedco2energylevelfigure}
  \lean{PhyXMiniProblems.ProblemPhyXMini0579.MatchesSuppliedCO2EnergyLevelFigure}
  \uses{def:physics:phyx-mini-0579:phyxminiproblems-problemphyxmini0579-energyinelectronvolts, def:physics:phyx-mini-0579:phyxminiproblems-problemphyxmini0579-expectedarrowinitiallevel, def:physics:phyx-mini-0579:phyxminiproblems-problemphyxmini0579-expectedarrowfinallevel, def:physics:phyx-mini-0579:phyxminiproblems-problemphyxmini0579-expectedarrowdirection, def:physics:phyx-mini-0579:phyxminiproblems-problemphyxmini0579-marsco2lasersetup}
  Primary-image evidence: all three level lines and both arrows are present; the level labels are 0, 0.165, and 0.289 eV; and the physical level energies agree with those printed readouts
\end{definition}
\begin{proof}
  This is the declared model definition of Matches Supplied CO2 Energy Level Figure.
\end{proof}
\begin{definition}[Has Physical CO2 Laser Quantities]
  \label{def:physics:phyx-mini-0579:phyxminiproblems-problemphyxmini0579-hasphysicalco2laserquantities}
  \lean{PhyXMiniProblems.ProblemPhyXMini0579.HasPhysicalCO2LaserQuantities}
  \uses{def:physics:phyx-mini-0579:phyxminiproblems-problemphyxmini0579-lengthinmeters, def:physics:phyx-mini-0579:phyxminiproblems-problemphyxmini0579-energyinjoules, def:physics:phyx-mini-0579:phyxminiproblems-problemphyxmini0579-frequencyinhertz, def:physics:phyx-mini-0579:phyxminiproblems-problemphyxmini0579-planckconstantinjouleseconds, def:physics:phyx-mini-0579:phyxminiproblems-problemphyxmini0579-marsco2lasersetup}
  Positivity and ordering conditions for the physical quantities
\end{definition}
\begin{proof}
  This is the declared model definition of Has Physical CO2 Laser Quantities.
\end{proof}
\begin{definition}[Uses Ordinary Planck Constant]
  \label{def:physics:phyx-mini-0579:phyxminiproblems-problemphyxmini0579-usesordinaryplanckconstant}
  \lean{PhyXMiniProblems.ProblemPhyXMini0579.UsesOrdinaryPlanckConstant}
  \uses{def:physics:phyx-mini-0579:phyxminiproblems-problemphyxmini0579-planckconstantinjouleseconds, def:physics:phyx-mini-0579:phyxminiproblems-problemphyxmini0579-planckconstantsivalue, def:physics:phyx-mini-0579:phyxminiproblems-problemphyxmini0579-marsco2lasersetup}
  The exact SI calibration of ordinary Planck's constant. Physlib provides a scalar reduced Planck constant, but no dimensionful ordinary h suitable for this photon-wavelength model
\end{definition}
\begin{proof}
  This is the declared model definition of Uses Ordinary Planck Constant.
\end{proof}
\begin{definition}[Satisfies CO2 Laser Laws]
  \label{def:physics:phyx-mini-0579:phyxminiproblems-problemphyxmini0579-satisfiesco2laserlaws}
  \lean{PhyXMiniProblems.ProblemPhyXMini0579.SatisfiesCO2LaserLaws}
  \uses{def:physics:phyx-mini-0579:phyxminiproblems-problemphyxmini0579-lengthinmeters, def:physics:phyx-mini-0579:phyxminiproblems-problemphyxmini0579-energyinjoules, def:physics:phyx-mini-0579:phyxminiproblems-problemphyxmini0579-frequencyinhertz, def:physics:phyx-mini-0579:phyxminiproblems-problemphyxmini0579-planckconstantinjouleseconds, def:physics:phyx-mini-0579:phyxminiproblems-problemphyxmini0579-speedoflightinmeterspersecond, def:physics:phyx-mini-0579:phyxminiproblems-problemphyxmini0579-marsco2lasersetup}
  Governing laws for the E₂ → E₁ laser emission: the photon carries the molecular level-energy drop; the Planck--Einstein relation is E = h ν; the vacuum light-wave relation is c = λ ν. These are general modeling relations among independent observables. None contains 10.0 μm or an answer-choice label
\end{definition}
\begin{proof}
  This is the declared model definition of Satisfies CO2 Laser Laws.
\end{proof}
\begin{definition}[Answer Choice]
  \label{def:physics:phyx-mini-0579:phyxminiproblems-problemphyxmini0579-answerchoice}
  \lean{PhyXMiniProblems.ProblemPhyXMini0579.AnswerChoice}
  Labels of the four wavelength choices printed with the problem
\end{definition}
\begin{proof}
  This is the declared model definition of Answer Choice.
\end{proof}
\begin{definition}[displayed Wavelength In Micrometers]
  \label{def:physics:phyx-mini-0579:phyxminiproblems-problemphyxmini0579-displayedwavelengthinmicrometers}
  \lean{PhyXMiniProblems.ProblemPhyXMini0579.displayedWavelengthInMicrometers}
  \uses{def:physics:phyx-mini-0579:phyxminiproblems-problemphyxmini0579-answerchoice}
  Wavelength in micrometres printed beside each answer label
\end{definition}
\begin{proof}
  This is the declared model definition of displayed Wavelength In Micrometers.
\end{proof}
\begin{definition}[recorded Dataset Answer]
  \label{def:physics:phyx-mini-0579:phyxminiproblems-problemphyxmini0579-recordeddatasetanswer}
  \lean{PhyXMiniProblems.ProblemPhyXMini0579.recordedDatasetAnswer}
  \uses{def:physics:phyx-mini-0579:phyxminiproblems-problemphyxmini0579-answerchoice}
  The answer label recorded by the source dataset, retained as metadata
\end{definition}
\begin{proof}
  This is the declared model definition of recorded Dataset Answer.
\end{proof}
\begin{definition}[displayed Wavelength Tolerance In Micrometers]
  \label{def:physics:phyx-mini-0579:phyxminiproblems-problemphyxmini0579-displayedwavelengthtoleranceinmicrometers}
  \lean{PhyXMiniProblems.ProblemPhyXMini0579.displayedWavelengthToleranceInMicrometers}
  Half of the 0.1 μm precision displayed by choice C. This represents the rounding implicit in the requested multiple-choice wavelength
\end{definition}
\begin{proof}
  This is the declared model definition of displayed Wavelength Tolerance In Micrometers.
\end{proof}
\begin{definition}[Matches Displayed Wavelength]
  \label{def:physics:phyx-mini-0579:phyxminiproblems-problemphyxmini0579-matchesdisplayedwavelength}
  \lean{PhyXMiniProblems.ProblemPhyXMini0579.MatchesDisplayedWavelength}
  \uses{def:physics:phyx-mini-0579:phyxminiproblems-problemphyxmini0579-lengthinmicrometers, def:physics:phyx-mini-0579:phyxminiproblems-problemphyxmini0579-marsco2lasersetup, def:physics:phyx-mini-0579:phyxminiproblems-problemphyxmini0579-answerchoice, def:physics:phyx-mini-0579:phyxminiproblems-problemphyxmini0579-displayedwavelengthinmicrometers, def:physics:phyx-mini-0579:phyxminiproblems-problemphyxmini0579-displayedwavelengthtoleranceinmicrometers}
  The physical lasing wavelength rounds to the wavelength of a choice
\end{definition}
\begin{proof}
  This is the declared model definition of Matches Displayed Wavelength.
\end{proof}
\begin{definition}[Is Unique Matching Displayed Wavelength]
  \label{def:physics:phyx-mini-0579:phyxminiproblems-problemphyxmini0579-isuniquematchingdisplayedwavelength}
  \lean{PhyXMiniProblems.ProblemPhyXMini0579.IsUniqueMatchingDisplayedWavelength}
  \uses{def:physics:phyx-mini-0579:phyxminiproblems-problemphyxmini0579-marsco2lasersetup, def:physics:phyx-mini-0579:phyxminiproblems-problemphyxmini0579-answerchoice, def:physics:phyx-mini-0579:phyxminiproblems-problemphyxmini0579-matchesdisplayedwavelength}
  Exactly one printed choice agrees with the modeled wavelength
\end{definition}
\begin{proof}
  This is the declared model definition of Is Unique Matching Displayed Wavelength.
\end{proof}
\begin{lemma}[laser Photon Energy In Electron Volts eq]
  \label{lem:physics:phyx-mini-0579:phyxminiproblems-problemphyxmini0579-laserphotonenergyinelectronvolts-eq}
  \lean{PhyXMiniProblems.ProblemPhyXMini0579.laserPhotonEnergyInElectronVolts_eq}
  \uses{def:physics:phyx-mini-0579:phyxminiproblems-problemphyxmini0579-energyinelectronvolts, def:physics:phyx-mini-0579:phyxminiproblems-problemphyxmini0579-marsco2lasersetup, def:physics:phyx-mini-0579:phyxminiproblems-problemphyxmini0579-matchessuppliedco2energylevelfigure, def:physics:phyx-mini-0579:phyxminiproblems-problemphyxmini0579-satisfiesco2laserlaws}
  The E₂ → E₁ transition has energy gap 0.289 eV - 0.165 eV = 0.124 eV
\end{lemma}
\begin{proof}
  The conclusion follows from the governing relations and figure readouts named in the statement of laser Photon Energy In Electron Volts eq.
\end{proof}
\begin{lemma}[lasing Wavelength In Meters eq]
  \label{lem:physics:phyx-mini-0579:phyxminiproblems-problemphyxmini0579-lasingwavelengthinmeters-eq}
  \lean{PhyXMiniProblems.ProblemPhyXMini0579.lasingWavelengthInMeters_eq}
  \uses{def:physics:phyx-mini-0579:phyxminiproblems-problemphyxmini0579-lengthinmeters, def:physics:phyx-mini-0579:phyxminiproblems-problemphyxmini0579-electronvoltinjoules, def:physics:phyx-mini-0579:phyxminiproblems-problemphyxmini0579-planckconstantsivalue, def:physics:phyx-mini-0579:phyxminiproblems-problemphyxmini0579-speedoflightinmeterspersecond, def:physics:phyx-mini-0579:phyxminiproblems-problemphyxmini0579-marsco2lasersetup, def:physics:phyx-mini-0579:phyxminiproblems-problemphyxmini0579-matchessuppliedco2energylevelfigure, def:physics:phyx-mini-0579:phyxminiproblems-problemphyxmini0579-hasphysicalco2laserquantities, def:physics:phyx-mini-0579:phyxminiproblems-problemphyxmini0579-usesordinaryplanckconstant, def:physics:phyx-mini-0579:phyxminiproblems-problemphyxmini0579-satisfiesco2laserlaws}
  Combining the energy drop, E = h ν, and c = λ ν gives the exact coherent-SI expression for the lasing wavelength
\end{lemma}
\begin{proof}
  The conclusion follows from the governing relations and figure readouts named in the statement of lasing Wavelength In Meters eq.
\end{proof}
% --- Archon declaration topology end ---
% --- Archon physics formalization source end ---
